\subsection{The Center of a Quaternion Algebra}

Let \(\alpha, \beta, \gamma\) be elements of \(K\) , and let \(F\) be the quaternion algebra of type \((\alpha, \beta, \gamma)\) .

PROPOSITION 2. --- a) Suppose that the field K has characteristic different from 2. If \(\gamma\) or \(4\alpha + \beta^{2}\) is nonzero, then the center of F is equal to K; otherwise, it has dimension 2 and is generated by 1 and ij - ji.

\begin{enumerate}
\def\labelenumi{\alph{enumi})}
\setcounter{enumi}{1}
\tightlist
\item
  Suppose that the field \(\mathrm{K}\) has characteristic 2. If \(\beta \neq 0\) , then the center of \(\mathrm{F}\) is equal to \(\mathrm{K}\) ; if \(\beta = 0\) , then the algebra \(\mathrm{F}\) is commutative.
\end{enumerate}

By formula (30) of III, §2, No.~5, p.~444, we have

\[
i j - j i = - \beta j + 2 k, \quad j k - k j = \beta \gamma - 2 \gamma i, \quad k i - i k = - 2 \alpha j - \beta k. \tag {4}
\]

An element \(q = x + yi + zj + tk\) of F is central if and only if it commutes with i and j, that is, we have

\[
2 z + \beta t = - \beta z + 2 \alpha t = 0 \quad \text { and } \quad 2 \gamma t = \beta \gamma t = 2 y = \beta y = 0.\tag{5}
\]

First suppose that the characteristic of K is different from 2. If \(\gamma\) is nonzero, then the equalities in (5) imply y = t = 0 and then z = 0; consequently, we have \(q \in K\) . If \(\gamma = 0\) and \(4\alpha + \beta^{2} \neq 0\) , then they imply y = z = t = 0, so we have \(q \in K\) . If \(\gamma = 4\alpha + \beta^{2} = 0\) , then the system (5) reduces to \(y = 2z + \beta t = 0\) , so we have \(q = x + t/2(ij - ji)\) , which completes the proof of a).

Now suppose that the field K has characteristic 2. The system (5) can then be written as \(\beta t = \beta z = \beta y = 0\) ; assertion b) follows.

